\section{Intellectual lineage, absorbed mechanisms and novelty boundary}
A central design principle of \RAKL{} is that a research methodology should learn from prior methods in the same way it learns from domain science: decompose external work into scoped mechanisms, retain provenance, normalize equivalent formulations, and narrow novelty whenever a prior mechanism is found. The bibliography is therefore not merely a list of neighbouring systems. It is an \emph{assimilation ledger}: each retained reference motivates, constrains or removes novelty from a specific RAKL operator.

\subsection{Automated science and self-improving agents}
Scientific-agent systems have already established literature-grounded hypothesis generation, multi-agent scientific workflows, long-horizon external research state and end-to-end automation. Co-Scientist combines generation, reflection, ranking and hypothesis evolution and includes experimental biomedical validation \cite{gottweis2026}; Robin integrates literature and data-analysis agents across a discovery workflow \cite{ghareeb2026}; AI Scientist-v2 automates experiment design, execution, analysis and paper generation using agentic tree search \cite{yamada2025}; Kosmos uses a structured world model for long-horizon discovery \cite{mitchener2025}; SciAgents couples multi-agent reasoning with ontological knowledge graphs \cite{ghafarollahi2024}; and PaperQA2 demonstrates strong agentic literature synthesis \cite{skarlinski2024}. Darwin G\"odel Machine, EvoSkill, SkillFoundry and EvoAgentBench further make recursive modification, failure-driven skill creation, held-out validation and ability transfer active prior-art areas \cite{zhang2025,alzubi2026,shen2026,gao2026}. \RAKL{} therefore does not claim novelty for ``AI scientist'', multi-agent debate, skill libraries or self-modifying code in isolation. Earlier LLM-agent work already demonstrates iterative self-feedback without protected scientific authority: Reflexion stores verbal feedback from trial-and-error in episodic memory \cite{shinn2023}, while Self-Refine uses the same model to generate, critique and revise outputs iteratively \cite{madaan2023}. These systems further narrow the RAKL claim: reflection, feedback memory and iterative refinement are prior mechanisms; the candidate contribution lies in subjecting method change to evidence lineage, negative-history retention and fresh assurance outside the challenger's write authority.

\subsection{From workflow completion to evidence-bounded scientific behavior}

The autonomous-science literature now makes a useful distinction between \emph{workflow coverage} and \emph{scientific reliability}. A system may be able to search, write code, call instruments, fit models and produce a manuscript while still failing at evidence attribution, uncertainty handling or revision after decisive counterevidence. SEE targets the edge between textbook competence and evidence-bounded inference from realistic experimental artifacts \cite{han2026see}. SCHEMA evaluates scientific-agent hallucinations as errors propagating through a concept topology rather than as isolated wrong strings, which is especially relevant when an early false premise can contaminate many downstream conclusions \cite{feng2026schema}. These evaluation directions are close to the \RAKL{} concern but differ in role. They are measurement frameworks for scientific-agent behavior; \RAKL{} proposes an internal state-transition discipline intended to make some of those failures representable and blockable.

Two newer benchmark families sharpen the fail-closed requirement. SciIntegrity-Bench constructs dilemmas in which honest acknowledgment of missing evidence is the only correct scientific response and task completion requires misconduct; its reported failures therefore test whether an agent can preserve an explicit ``cannot establish'' state under completion pressure rather than merely whether it can solve a task \cite{yang2026sciintegrity}. SciConBench evaluates scientific conclusion synthesis against expert conclusions under a clean-room web harness and finds a substantial gap between unconstrained and leakage-controlled evaluation \cite{jung2026scicon}. These results motivate two separate RAKL gates. Missing evidence must remain a typed block rather than be filled by synthetic observations, and evaluation evidence must be separated from information available to the proposal-generating process. Neither benchmark is evidence that RAKL passes those gates in the wild; each is prior art for a failure class the framework claims to make auditable.

AgentBuild offers a complementary engineering precedent. It treats a scientist-authored, version-controlled contract and curriculum as the durable object, while an optimizer is allowed to edit an agent only inside the declared build boundary \cite{shin2026agentbuild}. This is consonant with the \RAKL{} separation between replaceable proposal machinery and protected scientific or evaluative state. The important novelty boundary is therefore not ``agents should have rubrics'' or ``agents should use tests.'' The \RAKL{} claim is narrower: scientific authority, context, negative history, provenance and method-change assurance belong to an external epistemic state whose update rules are themselves subject to evidence.

These distinctions suggest a useful evaluation matrix. One axis asks \emph{how much of the scientific workflow can the system execute?} A second asks \emph{how often are scientific state transitions licensed correctly?} A third asks \emph{how well does the system discover actions that reduce the scientifically relevant uncertainty?} A fourth asks \emph{what resources are required?} A system can improve on one axis while regressing on another. For example, a broader tool set may increase workflow completion while increasing unsupported integration of tool outputs; an aggressive context compressor may reduce cost while dropping a falsifier; a self-editing agent may improve development scores while overfitting the evaluator. A single aggregate score can conceal these tradeoffs, which is why the present evaluation model retains blocking coordinates.

\subsection{Model building, prediction and mechanism}
Statistical methodology has long treated model construction as iterative rather than final. Box emphasized model criticism and revision as part of scientific modelling \cite{box1976}; posterior predictive assessment operationalizes comparison between observed discrepancies and model-generated replications \cite{gelman1996}. Breiman contrasted stochastic data models with algorithmic predictive modelling \cite{breiman2001}, while Shmueli made the distinction between explanatory and predictive goals explicit \cite{shmueli2010}. Mechanistic philosophy separately analyzes mechanisms in terms of organized entities and activities that produce regular change \cite{machamer2000}. These lines motivate RAKL's refusal to let predictive success mint mechanism authority. The new model-criticism operator is not claimed as a new statistical idea; RAKL's contribution is to make adequacy a typed child operation inside an evidence-governed research state. Popper's falsifiability and corroboration programme is also classical prior art for designing claims that can fail \cite{popper1959}. More recently, automated symbolic discovery has shown that mathematical regularities can be searched directly from experimental data \cite{schmidt2009}; RAKL treats such equation search as a proposal mechanism whose variables, observation maps, limiting cases and residuals still require scientific interpretation and verification.

Causal inference supplies a separate ancestry for intervention and mechanism language. Structural causal models distinguish observed association from intervention and counterfactual questions \cite{pearl2009}; modern causal-inference treatments likewise make assumptions and intervention distributions explicit \cite{peters2017}. RAKL absorbs this separation rather than using ``mechanism'' as a synonym for prediction. A predictive survivor may nominate a causal model, but causal or mechanistic authority requires additional assumptions and evidence.

\subsection{Prediction, mechanism, identification and explanation are different targets}

The distinction between prediction and mechanism has become more important as scientific AI systems improve at interpolation and forecasting. Mechanistic World Models argue that autonomous discovery requires reusable explanatory structure rather than predictive mappings alone \cite{posner2026mwm}. CausaLab operationalizes a related separation by evaluating both performance on tasks and recovery of a hidden causal mechanism, showing that the two can diverge substantially \cite{yang2026causalab}. In mechanistic interpretability, Lin and Liu make the parallel point that validation evidence such as ablations, faithfulness or completeness does not by itself identify a causal mechanism unless the required identification assumptions are stated \cite{linliu2026identification}. These works narrow rather than enlarge the \RAKL{} novelty claim: the representation--mechanism--identification separation is not a newly discovered philosophical fact.

\RAKL{} turns that separation into transition permissions. Let \(r\) be a representation that predicts an observable, \(m\) a candidate generating mechanism, and \(i\) an identification certificate. Evidence can increase representation authority without increasing mechanism authority:
\[
\Delta R(r)>0 \not\Rightarrow \Delta M(m)>0.
\]
Likewise, evidence that a mechanism class is scientifically plausible does not imply that the observations point-identify one member:
\[
M(m)>0 \not\Rightarrow I(m)=\text{POINT\_IDENTIFIED}.
\]
The correct terminal object may be an equivalence class, an identified set, a bound or a decision rule robust across several mechanisms. Manski's recent treatment of inductive risk under underdetermined theory choice reinforces why partial identification is not merely a statistical fallback: decision-making can remain principled when the evidence does not justify one uniquely selected theory \cite{manski2026underdetermination}.

This matters for scientific-agent design because generative systems are naturally pressured toward a singular narrative. A language model is rewarded for producing one coherent explanation, a model-selection pipeline returns a winner, and a paper conventionally tells one story. Those interface pressures can create an \emph{identification illusion}: the system mistakes a selected representative for an identified generator. \RAKL{} instead stores the survivor set and the observation map that produced the equivalence. If two mechanisms produce the same registered observations, selecting between them requires a discriminator that changes the observation boundary. If no feasible discriminator exists, the system should preserve the ambiguity and move downstream decisions to robust or set-valued analysis.

The distinction also clarifies mechanism ancestry. A mechanism claim is not strengthened merely by adding more descriptive detail. It requires a trace from lower-level objects and interactions through intermediate state and the observation process. The ancestry can be coarse and domain-specific; \RAKL{} does not require microscopic reductionism. What it forbids is a missing explanatory bridge being silently supplied by linguistic plausibility. Where the bridge is unavailable, the model is labelled phenomenological, predictive, teacher-only or otherwise scoped rather than being promoted by prose.

\subsection{Belief maintenance, negative history and partial identification}
Truth-maintenance systems already provide dependency-aware belief revision and inconsistency handling \cite{doyle1979}; assumption-based TMS extends this to multiple assumption-labelled environments and competing solutions \cite{dekleer1986}. AGM belief revision gives a separate axiomatic tradition for contraction and revision under minimal change \cite{agm1985}. RAKL therefore does not claim the generic idea of belief revision. Its negative-history rule is a scientific-governance specialization: a null, refutation or failed transition remains an addressable historical event even when later evidence supersedes its scope. Likewise, partial identification is established statistical methodology: Manski's bounds exemplify learning identified sets under weak assumptions \cite{manski1990}, while modern sensitivity analysis quantifies how conclusions change under plausible omitted-assumption worlds \cite{cinelli2020}. RAKL treats `PARTIALLY\_IDENTIFIED' and `ASSUMPTION\_SENSITIVE' as legitimate scientific outcomes rather than forcing point answers.

\subsection{Belief revision, argumentation and the status of negative evidence}

The belief-maintenance lineage explains why revision cannot be implemented as destructive overwrite. Doyle's truth-maintenance system records justifications so downstream beliefs can be withdrawn when their supports fail, while assumption-based TMS preserves multiple environments rather than forcing a single consistent context \cite{doyle1979,dekleer1986}. AGM revision supplies a complementary axiomatic view of minimal change under incorporation of new information \cite{agm1985}. Abstract argumentation adds a different object: a network in which arguments can attack one another and acceptability depends on the chosen semantics \cite{dung1995}. These frameworks are not interchangeable. TMS emphasizes dependency maintenance, AGM emphasizes rational belief-set change, and argumentation emphasizes dialectical support/attack structure. \RAKL{} borrows the need to preserve dependencies and incompatibilities, but its canonical objects are scientific claims scoped by evidence, context and authority rather than unqualified propositions.

This difference is most visible for negative evidence. A null, refutation or failed mapping has two temporal identities: it is evidence about the historical state of inquiry, and it may or may not remain decisive for a later, differently scoped claim. Deleting the old object when a successor arrives destroys information needed to audit cherry-picking, repeated testing, route resurrection and model-family evolution. The rule
\[
\mathcal H^-_t\subseteq \mathcal H^-_{t+1}
\]
therefore means preservation of the event and its original scope, not that every old negative conclusion remains eternally applicable. A later observation can narrow, supersede or contextualize the interpretation while leaving the receipt addressable.

This is also a response to a known scientific-selection problem. The ``file drawer'' literature shows why invisibility of null results can distort the apparent evidence landscape \cite{rosenthal1979,franco2014}. In \RAKL{}, negative history is not merely a publication-policy preference. It is a state variable used by search, saturation and self-evolution. If a candidate mechanism has already failed under a particular regime, a later generator must either explain why the new context is materially different or inherit that negative evidence. If a method repair was tried and failed, the next Self-\RAKL{} generation cannot erase the failure to make its lineage look monotone.

Model pluralism further motivates resisting forced synthesis. Veit argues that many scientific purposes require sets of models whose functions depend on context and history rather than one canonical model \cite{veit2020pluralism}. \RAKL{} gives that pluralism an operational destination: a plural atlas can be a successful outcome when different local models remain useful under distinct contexts, while an identified set is appropriate when the models compete on the same target but observations cannot distinguish them. The two should not be confused. Plurality can arise from legitimate perspectival specialization; non-identification arises from insufficient discriminatory evidence.

\subsection{Scientific revision as an auditable transaction}

The preceding traditions suggest a stronger implementation rule: scientific revision should be represented as a transaction with explicit preconditions, postconditions and preserved supersession lineage. A conventional belief store can replace proposition \(p\) with \(\neg p\). A scientific atlas must additionally answer which population, measurement, assumptions and evidence made each version applicable. The transaction object is therefore richer than a truth-value flip. One useful abstract form is
\[
\rho=(c_{\mathrm{old}},c_{\mathrm{new}},e,\gamma,\tau,\omega),
\]
where \(e\) is the evidence that triggered revision, \(\gamma\) its context, \(\tau\) the typed relation between old and new claims, and \(\omega\) the authority effect that the evidence is licensed to create. The old claim remains historically addressable even when its current applicability becomes empty.

This representation distinguishes at least five revision patterns. \emph{Refutation} records evidence incompatible with a claim under aligned context. \emph{Scope restriction} preserves the claim on a narrower domain after a counterexample appears elsewhere. \emph{Supersession} replaces an effective or approximate law with a successor while preserving the regime in which the old law remains adequate. \emph{Decomposition} splits one coarse claim into several context-specific claims after hidden heterogeneity is discovered. \emph{Authority downgrade} retains the descriptive content but withdraws a stronger mechanism, identification or decision certificate when its prerequisite fails. These outcomes are scientifically different even if a final prose summary would say simply that ``the model was updated.''

Dependency propagation is also typed. If claim \(b\) was derived from \(a\), refuting \(a\) may invalidate \(b\). If \(b\) was independently measured and merely consistent with \(a\), the same refutation need not remove \(b\). A provenance graph that records only citation ancestry cannot always make this distinction; the transition relation must say whether an edge is justificatory, derivational, contextual, analogical or merely navigational. This is where truth-maintenance ideas meet epistemic provenance: dependencies matter because they determine which descendants should be reconsidered when evidence changes \cite{doyle1979,dekleer1986}.

Argumentation adds another useful boundary. An attack edge can express that two claims cannot both retain a particular authority under a registered context, but an attack is not itself a refutation. A methodological critique may attack an inference while leaving the measurement intact. A new experiment may attack one mechanism but leave a predictive relation untouched. A conflict between two observational studies may be suspended until population and measurement alignment is established. The atlas therefore does not ask only ``which argument wins?''; it asks which coordinates of which scientific objects remain licensed after the conflict is typed \cite{dung1995}.

Revision has a chronology obligation. The system must know whether a hypothesis, threshold, covariate set or interpretation was fixed before the evidence that appears to support it. Otherwise the same data can play both proposal and confirmation roles without disclosure. This is why negative history, preregistration and exact candidate identity belong to the same control architecture. They are different ways of preserving the temporal direction of scientific justification.

The transaction view also explains why a paper should not silently rewrite its own past. If a manuscript once claimed that the generic atom--witness structure was a mathematical lattice and later analysis showed that only a compatibility complex had been established, the correct repair is not merely a global search-and-replace. The project should preserve the terminology error as method history, document the narrower structure, and state which downstream claims changed. The current compatibility-complex repair follows that pattern. Scientific self-correction is more informative when the reader can see what was corrected and why.

Finally, auditable revision provides a natural interface to sensitivity analysis. Instead of asking whether one final conclusion is ``robust'', the system can replay the transition under alternative registered assumptions or context mappings and observe which authority coordinates change. A conclusion that survives many plausible revisions has a different scientific profile from one that depends on a single fragile alignment, even if both currently have the same point estimate. RAKL treats that profile as structured evidence rather than compressing it into a confidence adjective.

\subsection{Why evidential selection pressure belongs inside the model of research}

Scientific literature is not an unbiased sample of all completed inquiries. Selective publication, flexible analysis and repeated exploration can change the distribution of claims that enter a search system \cite{rosenthal1979,franco2014,ioannidis2005}. A literature agent that retrieves published claims perfectly can therefore inherit upstream selection bias perfectly. Evidence governance begins before language-model reasoning.

\RAKL{} does not attempt a universal correction for publication bias. Instead it records relevant selection information when available and avoids treating publication count as independent confirmation. Pre-registration status, registered reports, availability of null results, dataset ancestry, replication status and analysis chronology can all alter how a claim should be interpreted. These become context/evidence attributes rather than one global quality label.

The same reasoning applies to benchmark ecosystems. Scientific-agent benchmarks are published objects selected by designers, and method developers adapt to visible tasks. A result on a saturated public benchmark can overstate transfer if the benchmark has become part of the development environment. Fresh-assurance tasks and frozen evaluation chronology are intended to protect against this adaptive exposure at the method level.

This is another reason negative history has value beyond transparency. Failed or null results change the selection model of what has been tried. Preserving them can prevent repeated rediscovery of a dead end and can reveal that an apparently dominant positive literature sits on top of many unreported or locally stored failures. In a research organization, the internal atlas can therefore contain scientifically valuable evidence that never became a conventional publication.

The authority system treats these factors as modifiers and blockers rather than as automatic verdicts. A preregistered study can be badly measured; a post hoc analysis can reveal a real phenomenon; a highly replicated result can still fail to transport. The purpose is to preserve the variables needed for scientific judgment and downstream sensitivity analysis.

\subsection{Provenance and machine-readable research objects}
Database provenance work had already made the origin and derivation of view data a first-class problem, including the distinction between ``why'' and ``where'' provenance \cite{buneman2001}. The W3C PROV family formalizes entities, activities, agents and derivations as provenance objects \cite{moreau2013}; FAIR articulates findability, accessibility, interoperability and reuse as stewardship goals \cite{wilkinson2016}; RO-Crate packages data, software and methods with machine-readable metadata and relationships \cite{soilandreyes2022}; and nanopublications attach provenance to granular scientific assertions \cite{kuhn2018}. Agent-workflow provenance and verifiable agent-native publishing are also active contemporary directions \cite{souza2025,dogah2026}. RAKL therefore does not claim generic provenance or atomic statement packaging. Its narrower addition is to bind provenance to context-scoped scientific authority, negative history, evidence-lineage-aware saturation and target-conditioned context compilation. Computational reproducibility is likewise an established discipline: reproducible-computation rules emphasize recording how each result was produced \cite{sandve2013}, while good-enough scientific-computing practice emphasizes data, software, collaboration and project organization \cite{wilson2017}. RAKL's receipts, environment attestation and artifact manifests are an implementation of this broader reproducibility requirement, not a novelty claim by themselves.

Recent agent-verification work also shows why pooled factual support is insufficient. ProvenanceGuard evaluates whether an atomic claim is supported by the particular source to which it is attributed, exposing cross-source conflation as a distinct error from ordinary factuality \cite{alvarez2026provenanceguard}. A broader survey of evidence tracing and execution provenance likewise separates retrieved evidence, tool outputs, memory, intermediate claims, actions and final answers rather than treating an agent trace as one undifferentiated rationale \cite{wang2026traces}. RAKL adopts the same granularity but adds a scientific-authority boundary: a correctly attributed source can still be weak, context-incompatible, derivative of the same evidence root, or irrelevant to a mechanism claim. Source ownership is therefore necessary for some updates but never sufficient for arbitrary authority escalation.

\paragraph{Provenance standards and the novelty boundary.}
Provenance representation itself is established infrastructure. W3C PROV-O, for example, supplies a general ontology for interchanging provenance across heterogeneous systems \cite{w3cprov2013}. \RAKL{} therefore does not claim novelty for provenance graphs or source ancestry. Its narrower contribution is to couple source-rehydratable provenance to context-scoped scientific objects, typed relation witnesses, multi-axis authority and gated canonical updates; provenance visibility alone neither establishes truth nor licenses a mechanism or identification claim.

\subsection{From document provenance to epistemic provenance}

Machine-readable provenance solves an important but limited problem: it records where an object came from and how it was transformed. It does not, by itself, determine what the object is allowed to support scientifically. The Open Research Knowledge Graph (ORKG), for example, moves scholarly communication toward structured, machine-actionable research contributions rather than document-only representation \cite{jaradeh2019orkg}. W3C PROV supplies interoperable entity--activity--agent derivations \cite{w3cprov2013}. Nanopublications and RO-Crate provide other granular or packaged representations \cite{kuhn2018,soilandreyes2022}. These are essential precedents for a scientific state that must be inspectable by software.

\RAKL{} adds a second layer that we call \emph{epistemic provenance}. For an asserted scientific relation, the system must retain not only the source chain but the context alignment, verification identity, relation type, error semantics and authority effect. A source can be perfectly provenance-traceable and still be weak evidence; two sources can be independently published but descend from the same primary dataset; a derived analysis can be reproducible while its measurement operator is mismatched to the claim. Thus
\[
\text{traceable provenance}\not\Rightarrow
\text{independent evidence}\not\Rightarrow
\text{sufficient authority}.
\]
The evidence-lineage graph addresses the middle implication. It attempts to represent shared ancestry so that ten derivative reports of one measurement do not count as ten independent flat searches or ten independent confirmations.

Cognitive provenance is separated again. A workspace item may causally influence a proposal because it was retrieved or emphasized, yet that influence does not make the item evidence for the proposal's scientific content. The two edges have different semantics:
\[
\Pi^{\mathrm{cog}}\neq \Pi^{\mathrm{evid}}.
\]
This distinction becomes important when using LLM traces, memory systems or interpretability tools. A representation being computationally load-bearing is evidence about the agent's cognition; it is not automatically evidence about the external phenomenon being studied. A paper that mixes those graphs can accidentally turn explanation of the model into explanation of nature.

The practical consequence is an ``authority firewall'' around transformations. Raw acquisition can create a source object; projection can create a context-scoped claim object; normalization can create an equivalence witness; an embedding can create a retrieval index; compression can create a derivative view; a workspace can create an access relation. None of those operations, on its own, increases scientific authority. Only a verified promotion operation can do so. This makes provenance a necessary substrate for authority without making provenance itself authority.

\subsection{Analogy and active inquiry}
Structure-mapping theory treats analogy as relational mapping rather than surface resemblance \cite{gentner1983}; constraint-satisfaction accounts jointly evaluate competing mappings \cite{holyoak1989}. This prior art motivates RAKL's separation of conservative `GLUE' from exploratory `JUMP': an analogy must state preserved and non-preserved structure and has discovery authority only until target evidence arrives. Experimental information itself predates modern active learning: Lindley formalized expected information from an experiment \cite{lindley1956}, and Bayesian experimental-design reviews systematized utility-based design under prior/posterior uncertainty \cite{chaloner1995}. Active data selection also has a long history: information-based design chooses observations by their expected information value relative to a model class \cite{mackay1992}. RAKL uses information gain when calibrated probabilistic models are defensible, and otherwise falls back to set shrinkage or worst-case mechanism separation rather than fabricating priors. Repeated adaptive evaluation is itself a validity problem, motivating the protected fresh-assurance reserve \cite{dwork2015}.

\subsection{Discovery search, experimental design and stopping are one control problem at different layers}

Literature search and physical experimentation are often treated as separate modules, but both can be written as selection under uncertainty. A literature query chooses an observation operator over the published record; an experiment chooses an observation operator over the scientific system. In each case the purpose is not simply to collect more items. It is to expose distinctions that matter for the registered target. Bayesian experimental design formalizes one version through expected utility \cite{lindley1956,chaloner1995}; active learning chooses measurements that are expected to improve a learner \cite{mackay1992}. \RAKL{} uses these ideas only when the probability model required by the utility is itself defensible.

Literature-based discovery supplies an important cross-domain precedent for the search side. Swanson's classic ``undiscovered public knowledge'' example showed that useful connections can exist between literatures that are individually public yet not directly connected \cite{swanson1986}. This is precisely why semantic saturation cannot be established by repeating the framework's own vocabulary. Open-World Mechanism Discovery therefore includes function-only, historical, mathematical, implementation, citation-neighborhood, adversarial and freshness routes. A later exogenous concept that changes the novelty boundary is a falsifier of the prior route-closure claim for that scope.

Stopping research introduces a dual risk. Stopping too early misses a material semantic object; stopping too late wastes resources and creates more opportunities for adaptive overfitting. Technology-assisted systematic review has developed explicit stopping methods because screening ``until tired'' is not a defensible rule. Recent confidence-based approaches shift attention from raw recall toward whether the information need is sufficiently resolved \cite{fletcher2026stopping}. \RAKL{} adopts the general lesson without importing a particular recall estimator: stopping must be attached to a registered information need, source/query universe, evidence lineage and reopen condition.

This produces a hierarchy of stopping statements. ``Budget exhausted'' is a resource state. ``Route locally flat'' means one query family yielded no new canonical semantic object. ``Same-context plateau'' means repeated passes in one research context are flat. ``Evidence-lineage independent flatness'' requires genuinely distinct evidence ancestry. ``Scoped saturation'' additionally requires route coverage, unresolved-obligation accounting and no native residual that reopens the fiber. None of these implies that the open world contains no unknown relevant concept.

\subsection{Metacognition and learning from challenge}
RAKL's human-learning analogy is functional rather than anthropomorphic. Self-regulated learning distinguishes forethought, monitoring/control and reflection \cite{zimmerman2002}; productive-failure work shows that initial failure can improve later representation and transfer under appropriate conditions \cite{kapur2008}. Psychology also documents reasons not to trust verbal introspection alone: people can underestimate their own bias \cite{pronin2002}, overestimate explanatory understanding until forced to reconstruct an explanation \cite{rozenblit2002}, and benefit in some settings from explicitly considering the opposite hypothesis rather than generic instructions to ``be unbiased'' \cite{lord1984}. These results motivate executable RAKL controls such as calibration audits, explanation reconstruction, countermodels, strategy switching, help-seeking, anti-rumination and fresh transfer. Recent generative-AI work similarly argues for metacognitive control around rather than merely inside the model \cite{ji2026}.

\subsection{Why reflection is not an evaluator}

A self-reflective LLM can notice errors, but the act of reflection is not statistically or epistemically independent of the process that generated those errors. This matters because many contemporary agent loops reuse one model as proposer, critic, judge and editor. Such role separation can be operationally useful while leaving shared blind spots, shared training priors and shared context contamination intact. The psychology literature on bias blind spots and illusion of explanatory depth provides human analogies for this failure \cite{pronin2002,rozenblit2002}; \RAKL{} treats the analogy as motivation, not as a claim that language models reproduce the same cognitive mechanisms.

The architecture therefore separates \emph{metacognitive diagnosis} from \emph{assurance}. Diagnosis asks what kind of failure occurred and what action should be taken next. Assurance asks whether a frozen candidate method change transfers on a task or realization that did not participate in designing the repair. The first can be same-context. The second must be protected to the degree required by the claim. Development improvement \(\Delta_D>0\) is thus necessary but insufficient for strong self-evolution; the protected comparison must also show \(\Delta_A>0\) under unchanged blocking criteria.

This separation also prevents ``reflection inflation.'' Adding more self-critique turns is not automatically more metacognition. A useful metacognitive operation must change a control decision: acquire evidence, switch strategy, construct a countermodel, ask for external help, stop an unproductive route, or open a method-basis gap. Repeating explanations without changing an uncertainty-bearing operation is recorded as flat effort. That is why learning-progress signals are coupled to action policy rather than treated as authority.

Finally, metacognitive state itself must remain revisable. A controller can misattribute a data problem to a method problem or a method problem to a model problem. Self-\RAKL{} therefore treats failure attribution as a hypothesis with discriminators. Only a supported method-basis gap justifies inventing or assimilating a new operator. Missing measurements route to evidence acquisition; implementation defects route to repair; objective defects route to governance. The ability to invent a method is not permission to invent one whenever performance is poor.

\subsection{Memory, retrieval and local-to-global knowledge}
Retrieval-augmented generation already combines parametric generation with explicit non-parametric memory \cite{lewis2020}; therefore external knowledge access is not itself a RAKL contribution. Hierarchical memory and context optimization are prior engineering domains. MemGPT frames long-horizon context as virtual memory \cite{packer2023}; RAPTOR organizes multi-resolution retrieval hierarchically \cite{sarthi2024}; LLMLingua compresses prompts \cite{jiang2023}; and RECOMP selectively compresses retrieved material \cite{xu2023}. Long-context models also do not necessarily use all prompt positions equally: ``lost in the middle'' evaluations show systematic degradation when relevant information lies away from context boundaries \cite{liu2024}. This supports treating active scientific context as an engineered working set rather than assuming that a maximal context window is equivalent to usable memory. Sheaf-theoretic work has long characterized contextuality through obstructions to global sections \cite{abramsky2011}; Local-to-global consistency and provenance-rich contextual worlds are also active formal directions \cite{gibson2026,vitali2026}. RAKL therefore does not claim generic graph storage, compression, hierarchical retrieval or sheaf mathematics. Its narrower requirement is that scientific projection, semantic identity, storage compression and prompt selection remain distinct operations, and that mandatory epistemic evidence cannot be silently traded away for context-window efficiency.

Taken together, these literatures narrow the candidate novelty to the integrated transition discipline: source contribution as contextual projection before competition; typed non-escalating authority; obstruction-aware local-to-global synthesis; immutable negative research history; evidence-lineage-aware stopping; bounded but rehydratable scientific memory; residual-driven method completeness; and protected recursive method evolution.

\paragraph{Compression lineage and what is not claimed.}
Minimum-description-length reasoning formalizes an important pressure against representational complexity \cite{rissanen1978}, while the information-bottleneck principle studies compression that preserves information relevant to a target variable \cite{tishby2000}. These are intellectual-lineage analogies rather than algorithms silently implemented by the current reference profile. The present context compiler does not estimate mutual information or solve an information-bottleneck objective, and the Knowledge Atlas ``volume'' introduced below is not a description length. The retained engineering lesson is narrower: compactness, target-relevant retained function and scientific authority must be measured separately.

\subsection{Local-to-global mathematics: sheaves, closure systems and why the names matter}

Several mathematical traditions offer precise versions of ``local pieces forming a global object,'' and their differences help prevent metaphor from becoming false formalism. Sheaf theory studies compatible local data and the conditions under which local sections can be glued into global sections. In contextuality, the absence of a global section can itself carry substantive meaning \cite{abramsky2011}. Recent AI-science work uses sheaf-theoretic transport and obstruction to diagnose whether a representational language can coherently extend into a new regime \cite{olivieri2026sheaf}. This is close to \RAKL{} in vocabulary and motivation. It narrows the novelty claim and supplies a useful comparator: the cited framework is a finite diagnostic of transport and theory-shift candidates, while \RAKL{} embeds obstruction inside a larger evidence, authority, memory, search and method-evolution control plane.

Formal Concept Analysis (FCA) offers a different route to genuine lattice structure. Given a formal context relating objects to attributes, FCA constructs formal concepts and orders them into a concept lattice \cite{ganterwille2024}. This matters because the historical name ``Recursive Atomic Knowledge Lattices'' can be read as an order-theoretic claim. The current atom--witness--path implementation does not obtain a lattice merely because atoms are combinable. Its honest generic structure is a typed compatibility complex. A local order-theoretic lattice can be licensed only after a scoped order or closure operator is defined and the relevant meet/join obligations are established. FCA is therefore a useful mathematical neighbor and a terminology constraint, not a hidden implementation detail.

The distinction can be summarized as follows. A graph gives adjacency. A compatibility complex gives witnessed co-admissibility among typed objects and constructive paths. A poset gives an antisymmetric order. A lattice additionally supplies meets and joins for the relevant pairs. A closure system supplies a family of closed objects and can induce lattice structure under the appropriate conditions. A sheaf organizes local data, restrictions and gluing. These structures can coexist in one research architecture, but none is interchangeable with another. \RAKL{} uses different ones for different jobs: provenance is graph-like, authority is partially ordered, local chart synthesis is atlas/sheaf-like when the obligations are defined, and constructive candidate assembly currently uses a compatibility complex.

This layered view also changes what ``global'' means. A global model may fail to exist because local charts genuinely conflict; it may exist but be non-unique because the evidence underdetermines it; or it may be representable but scientifically unsupported. Global existence, uniqueness and authority are therefore three separate predicates. A successful gluing calculation cannot by itself increase evidence quality, and strong evidence for local charts cannot guarantee a unique global synthesis.
